% ---------- ABSTRACT (spec section 4) : page iii, <= 500 words ----------
\phantomsection\addcontentsline{toc}{chapter}{\protect\hspace*{0.5in}Abstract}
\frontheading{Abstract}

People in aerial disaster imagery are difficult to detect because they are
often only a few pixels wide, partly occluded, and surrounded by rubble, water,
or smoke. This thesis asks a practical question: which changes to a compact
real-time detector improve this case enough to justify their computational
cost? Starting from YOLO11m, four configurations are trained under the same
protocol on the C2A disaster-imagery benchmark. The additive ablation replaces
the baseline attention module with CBAM, adds a stride-4 P2 detection scale,
and then tests a bidirectional state-space neck. The P2 scale produces the
clearest architectural gain. It raises recall for targets below 8~px from
$0.743$ to $0.757$ and $\mathrm{AP}_{50}$ from $0.843$ to $0.853$, while adding
about $0.5$~million parameters and $1$~ms of latency. CBAM reduces the model by
about $1$~million parameters while preserving the measured accuracy. The
state-space neck adds $2.4$~million parameters and increases latency by a
factor of 2.8 without an accuracy gain, so it is rejected.

The second phase tests whether these conclusions survive a stricter protocol
and real imagery. A split audit finds that 98\% of the official C2A test scenes
share a background scene with training. Retraining on a scene-disjoint split
reduces $\mathrm{AP}_{50}$ by $1.6$ points but does not change the architectural
conclusion. Further improvement comes from training assignment rather than
another feature branch. Blending normalised Wasserstein similarity into the
task-aligned assigner raises sub-8-pixel recall in all three replicated seeds
(mean $+1.6\pm0.5$ points); in the matched seed-0 pair the gain is $2.08$
points with a 95\% bootstrap interval of $[1.77,2.38]$. Overall $F_2$ and
$\mathrm{AP}_{50}$ remain neutral, showing that the method reallocates recall
towards the smallest targets instead of improving every size band. A
frequency-gated architecture branch fails its paired test and is also rejected.

Finally, real footage from the author's drone flights, campus scenes, and
disaster news video is used to measure the synthetic-to-real gap. The
synthetic-trained model scores $0.335$ $\mathrm{AP}_{50}$ on the drone test;
supervised adaptation with a few hundred labelled real frames raises this to
$0.795$ while retaining $0.837$ on C2A. The much larger within-video gain on
disaster footage does not reproduce on a held-out event, which limits the
generality of that result. The recommended CBAM+P2 model has $19.6$~million
parameters, runs in $14.6$~ms per image on an RTX 4070 Ti SUPER, and reaches
$0.853$ $\mathrm{AP}_{50}$ and $0.844$ $F_2$ on the official C2A test set. The
thesis therefore contributes both a compact detector and a controlled
record of which apparently promising changes did, and did not, improve tiny
aerial-human detection.
